\section{Introduction}
\label{sec:introduction}

Thermal models connect laser processing conditions to the melting and solidification histories that influence material structure in laser powder bed fusion (LPBF). Their usefulness extends beyond estimating the dimensions of a melted track: interpreting a process--structure relationship also requires temperature gradients and the thermal path followed as the material cools. Resolving the coupled laser, phase-change and transport processes over practical build scales remains computationally demanding~\cite{king2015}. This creates a persistent tension between affordable heat-transfer calculations, often constrained by accessible melt-pool geometry, and the credibility of the thermal interpretation drawn from them. Thermographic measurements provide additional information on spatial profiles and material histories~\cite{hooper2018,myers2023}, but geometric agreement remains an important practical constraint. Its meaning depends jointly on the heat-source representation, material response and definition of the measured observable.

The available modeling approaches address different parts of this tension. Powder-resolved thermo-fluid simulations describe recoil pressure, surface-tension-driven flow and evaporation in 316L, while ray-tracing/free-surface calculations have linked multiple reflections to keyhole-induced porosity in Ti6Al4V~\cite{khairallah2016,bayat2019}. Their physical detail is valuable for studying mechanisms that a conduction model represents only indirectly. Moving-source conduction provides a more economical thermal description. Rosenthal's solutions established heat-flow and cooling relations for moving sources; later finite-source and nonlinear formulations incorporated temperature-dependent properties and enthalpy storage~\cite{rosenthal1946,vanelsen2007}. Scaling relations based on normalized energy input and dwell-to-diffusion time further reduce the calculation needed to estimate dimensions~\cite{rubenchik2018}. This hierarchy permits useful simplification, provided the assessment addresses the output and operating regime for which the model is used.

Conduction-based models have achieved substantial geometric agreement under defined conditions. Promoppatum et al. found that both finite-element and Rosenthal approaches reasonably estimated IN718 melt-pool shapes, with differences in absorption sensitivity and high-input behavior~\cite{promoppatum2017}. For the AMMT IN625 single-track benchmark, Kollmannsberger et al. combined nonlinear phase-change conduction, measured boundary input and calibrated effective directional conductivity to reproduce the pool dimensions~\cite{kollmannsberger2019}. That result supports the utility of an effective thermal description; it also shows that source and directional material choices can be consequential when a scalar coupling factor alone cannot reproduce the shape. Across materials and regimes, Naderi et al. found that the fidelity of existing depth scaling laws varied and improved with revised, regime-specific relations~\cite{naderi2023}. Neither successful fitting nor a discrepancy therefore establishes a universal conclusion about conduction models.

Heat-source calibration has extended their practical range. Ross et al. inferred double-ellipsoid parameters from temperatures at the solidification boundary and obtained trends that permit interpolation to intermediate power and speed~\cite{ross2022}. Whalen et al. incorporated temperature-dependent and powder-bed properties in an adapted Eagar--Tsai model for 316L, inferred distributions of apparent absorption and porosity, and propagated geometric uncertainty~\cite{whalen2021}. More recently, Kusano and Watanabe compared four source families for IN738LC and used Bayesian optimization and regression to obtain a conical source applicable within their studied process window~\cite{heatsourcebayes2024}. Coleman et al. calibrated a dynamic volumetric source against bare-plate IN625 tracks at 195~W and 800~mm/s over multiple spot sizes, then applied it to an AMB2018-01 layer case~\cite{dynamic2024}. Their formulation includes temperature-dependent, phase-specific properties and latent storage. These achievements make source representation a substantive alternative to material closure when interpreting shape discrepancies.

Additional measurements expose the distinct information carried by geometry, energy coupling and temperature. Myers et al. demonstrated that parameter combinations fitting melt-pool geometry could produce different thermal profiles; two-color measurements helped distinguish those combinations~\cite{myers2023}. Calorimetric studies on bare and powder-coated materials show that effective absorption varies with power and surface condition~\cite{trapp2017}. For IN625, Lane et al. distinguished unreflected laser power from substrate heating because plume interactions and mass-energy losses also enter the energy accounting~\cite{lane2020absorption}. A fitted source factor consequently carries the meaning of the model and observation used to obtain it. Sensor-to-field mapping also matters: Weeks et al. modeled reflections of thermal emission within concave 316L keyholes to reconcile simulated and camera temperatures~\cite{weeks2026}. That observation correction addresses a different source of discrepancy from changing the underlying heat-transfer field.

Recent reliability studies make observable-specific assessment explicit. Soares et al. compared an analytical Gaussian conduction model with dual-wavelength imaging and metallography for 316L, IN625 and CoCr single tracks~\cite{soares2026}. With fixed properties selected near 80\% of the melting temperature and no explicit phase-change treatment, their model gave useful surface comparisons while depth correlation coexisted with substantial scale bias. The result distinguishes correlation across conditions from agreement in a particular dimension. The AMMT measurements likewise define imaging length and metallographic width/depth through different procedures and populations~\cite{lane2020}. A geometric constraint is therefore informative only in relation to its observation operator; fitting one dimension does not itself test the remaining shape or a material cooling history.

Material sensitivity and conduction-failure diagnosis are also established research directions. Yang et al. quantified the response of Ti6Al4V thermal fields and dimensions to processing and material parameters using constant-property analytical source models~\cite{yang2024}. Hong et al. examined 231 bare-plate tracks with an Eagar--Tsai conduction description, using measured coupling for IN718 and a published absorption relation for IN625 and 316L~\cite{hong2026}. Their inversion asks whether measured depth requires physically infeasible absorption, with effective-diffusivity contrasts testing the diagnosis. These studies address parameter sensitivity and model feasibility directly. A complementary need in a nonlinear, geometrically calibrated model is to resolve what a specified closure change does to the local phase-boundary structure, rather than infer that structure from a total dimension or a joint-property response.

This distinction is pertinent when material tables end below the temperatures at which phase change is represented. The IN625 supplier conductivity and sensible-heat-capacity inputs used here terminate below the adopted solidus~\cite{specialmetals625}. Continuing those tables introduces assumptions affecting both heat redistribution and storage. In an enthalpy formulation, sensible heat capacity and latent-inclusive storage are different quantities; changing both continuations together can conceal their separate responses. Similarly, movement of the rear solidus and liquidus positions together changes pool length without requiring comparable growth of the interval between them. That spatial phase interval and the time spent traversing it are also distinct, because material passage through a quasi-steady translating field depends on scan speed. Separating these responses makes the thermal meaning of a geometric sensitivity more explicit.

Here, a nonlinear moving-frame enthalpy model examines these distinctions for three AMMT bare-plate IN625 conditions. One effective source factor is fitted to the B surface-length proxy. Width and depth are excluded from that fit, while A and C provide retrospective, nonblind comparisons with the factor fixed. At B, a four-corner contrast separately replaces linear continuation by last-value holding for conductivity and sensible heat capacity, retaining the source factor and phase law without recalibrating the branches. The analysis separates their finite effects and interaction, decomposes rear solidus/liquidus displacement from phase-span change, and maps that span to passage time and interval-mean cooling rate using scan speed. Numerical verification supports these model comparisons; the derived histories remain distinct from independent thermal validation. The study thus provides a controlled diagnosis of thermal interpretation after one geometric calibration, with explicit distinctions between fusion-shape discrepancy, closure response and material time.
